\section{План реализации Fourier Neural Operator (FNO)}

\subsection{Постановка задачи}

Пусть требуется аппроксимировать нелинейный оператор
\[
\mathcal{G} : u(x) \rightarrow v(x),
\]
где $u(x)$ — входная функция (например, коэффициенты или начальные условия PDE),
а $v(x)$ — решение уравнения.

FNO обучает отображение между функциональными пространствами.

\subsection{Общая архитектура}

Архитектура FNO состоит из трёх этапов:

\begin{enumerate}
    \item Подъём размерности (Lift)
    \item Несколько Fourier-блоков
    \item Проекция в выходное пространство
\end{enumerate}

\subsection{1. Lift-слой}

Вход:
\[
u(x) \in \mathbb{R}^{N_x \times C_{\text{in}}}
\]

Применяется точечное линейное преобразование:
\[
w(x) = P u(x),
\]
где $P$ — обучаемая матрица.

Размерность:
\[
C_{\text{in}} \rightarrow C_{\text{width}}.
\]

\subsection{2. Fourier-слой}

Один FNO-блок реализует преобразование:

\[
w_{l+1}(x) =
\mathcal{F}^{-1} \left(
R(k)\,\mathcal{F}(w_l(x))
\right)
+
W w_l(x),
\]

где:

\begin{itemize}
    \item $\mathcal{F}$ — преобразование Фурье,
    \item $R(k)$ — обучаемые комплексные веса для первых $m$ мод,
    \item $W$ — точечное линейное преобразование,
    \item применяется нелинейность $\sigma$ (например GELU).
\end{itemize}

Полная формула слоя:

\[
w_{l+1}(x) =
\sigma\left(
\mathcal{F}^{-1} \left(
R(k)\,\mathcal{F}(w_l)
\right)
+
W w_l
\right).
\]

\subsection{Ограничение числа мод}

Используются только низкие частоты:

\[
R(k) = 0, \quad |k| > m.
\]

Это обеспечивает:

\begin{itemize}
    \item устойчивость
    \item снижение числа параметров
    \item регуляризацию модели
\end{itemize}

\subsection{3. Проекция}

После нескольких Fourier-блоков применяется:

\[
v(x) = Q w_L(x),
\]

где $Q$ — линейная проекция в выходное пространство
\[
C_{\text{width}} \rightarrow C_{\text{out}}.
\]

\subsection{Полная схема}

\[
u(x)
\rightarrow
\text{Lift}
\rightarrow
\text{FNO Block} \times L
\rightarrow
\text{Projection}
\rightarrow
v(x)
\]

\subsection{Алгоритм прямого прохода}

\begin{enumerate}
    \item Применить Lift-слой
    \item Для каждого Fourier-блока:
    \begin{enumerate}
        \item Выполнить FFT
        \item Умножить низкие моды на $R(k)$
        \item Выполнить обратное FFT
        \item Добавить линейную часть
        \item Применить нелинейность
    \end{enumerate}
    \item Применить выходную проекцию
\end{enumerate}

\subsection{Функция потерь}

Обычно используется среднеквадратичная ошибка:

\[
\mathcal{L}
=
\frac{1}{N}
\sum_i
\| v_{\text{pred}}(x_i) - v_{\text{true}}(x_i) \|^2.
\]
